\documentclass [11pt,a4paper,oneside,openany]{book} %Classe del documento, formato carta, singola facciata, apertura capitoli pag destra/sinista (ininfluente se impostato oneside). Formato libro
\usepackage[italian]{babel} %Lingua Documento (da impostare per la sillabazione)
\usepackage{graphicx} %Pacchetto necessario per la gestione delle immagini
\usepackage{xcolor,colortbl} %Colori
\usepackage{enumitem} % pacchetto per gestire le enumerazioni e resume
\usepackage {fancyhdr} %Pacchetto per la gestione accurata della pagina
\usepackage {setspace} %Pacchetto necessario per i comandi successivi onehalfspacing, singlespacing....
\usepackage{titlesec} % Pacchetto per le sezioni e sottosezioni
\usepackage{longtable} %Pacchetto per la gestione delle tabelle grandi
\usepackage[colorlinks=true]{hyperref} %Segnalibri nel pdf finale, sotto relativi parametri
\usepackage[T1]{fontenc}
\usepackage[utf8]{inputenc} %Pacchetti per le lettere accentate utf8
\usepackage{listings} %Pacchetto per formattare codice sorgente
\usepackage{mathtools} %Pacchetto matematica
\usepackage{float} %Pacchetto per il posizionamento delle figure
\usepackage{amsthm}


\hypersetup{
	bookmarksnumbered=true,
	linkcolor=black,
	citecolor=black,
	urlcolor=black,
}

\usepackage{geometry} % Dimensione pagina

\geometry{a4paper} % Formato carta
\addtolength{\textheight}{75pt} %Margini
\oddsidemargin 30pt

\setcounter{secnumdepth}{2}
\setcounter{tocdepth}{3}


\newcommand{\reqref}[1]{\textbf{R\ref{#1}}}
\newcommand{\goalref}[1]{\textbf{G\ref{#1}}}

\definecolor{darkgreen}{rgb}{0,0.6,0}
\lstset{frame=tb,
  language=Java,
  morekeywords={module, open, integer, ordering, as, sig, one, some, set, not, in, all, no, none, and, or, fact, pred, fun, implies, sum, check, run, but},
  aboveskip=3mm,
  belowskip=3mm,
  showstringspaces=false,
  columns=flexible,
  basicstyle={\small\ttfamily},
  numbers=none,
  numberstyle=\tiny\color{gray},
  keywordstyle=\color{blue},
  commentstyle=\color{darkgreen},
  %stringstyle=\color{mauve},
  breaklines=true,
  breakatwhitespace=true,
  tabsize=3
}

\begin{document} %Inizio Documento

%%FRONTESPIZIO%%=========================================
\begin{titlepage}

 \begin{center} 
     \includegraphics[height=6cm]{logo.eps}\\
     \vspace{4em}
     {\Large \textsc{Software Engineering II}}\\
     \vspace{6em}
     {\LARGE \textbf{Travlendar+}}\\
     \vspace{3em}
     {\Large \textsc{Requirements Analysis}}\\
     \vspace{1em}
     {\Large \textsc{and}}\\
     \vspace{1em}
     {\Large \textsc{Specifications Document}}\\
 \end{center}
 
    \vskip 2cm
 
 	Authors:
 	\vspace{0.5em}
 	\begin{center}
      {\Large \textit{Edoardo D'Amico}}\\
      {\Large \textit{Gabbolini Giovanni}}\\
      {\Large \textit{Parroni Federico}}\\
    \end{center}

\vskip 2.0cm
\begin{center}
{\normalsize 1$^{st}$ October 2017}
\end{center}

\end{titlepage}

\newpage

%%INTESTAZIONI PAGINA%%==================================
\pagestyle{fancy}
\renewcommand{\chaptername}{Chapter}
\renewcommand{\chaptermark} [1]{\chaptername\ \thechapter.\ #1}{} 
\renewcommand{\chaptermark}[1]{\markboth{\thechapter.\ #1}{}} 
\renewcommand{\sectionmark}[1]{\markright{\thesection\ #1}}
\fancyhf{}
\fancyhead[LE,RO]{\bfseries\thepage} 
\fancyhead[LO,RE]{\bfseries\leftmark}
\fancypagestyle{plain}{%
\fancyhead{} % leva l'intestazione
\renewcommand{\headrulewidth}{0pt} % e la linea
}


%%INDICE%%==============================================
\renewcommand{\contentsname}{Indice}
\renewcommand{\listfigurename}{Elenco delle Figure}
\renewcommand{\listtablename}{Elenco delle Tabelle}
\tableofcontents

\linespread{1.5}\selectfont		%interlinea
%%CAPITOLI===============================================
%%CAPITOLO 1: Introduzione=======================================
\chapter{Introduction}

\section{Purpose}
The purpose of this document is to provide more technical details than the RASD about Travelendar+ system. This document is addressed to developers and aims to identify:
\begin{itemize}
\item The high level architecture
\item The design patterns
\item The main components and their interfaces provided one for another
\item The Runtime behavior
\end{itemize}

\section{Scope}

Travlendar+ is an application that aims to ease the management of the daily appointments of a registered user by providing a desired schedule for those. 
The application runs within the context of a restricted geographic area and thus should consider only the travel means available in this scope.
The application is able to interact with different data sources, in particular Travel Means APIs and User device GPS, and to arrange the appointments by exploiting these information.
Once the schedule is computed, the application gives to the user the possibility to buy tickets for the involved travel means, then he/she can run it and follow directions provided by the application. 

\section{Definitions, Acronyms, Synonims}

\subsection{Definitions}
\theoremstyle{definition}
\newtheorem{definition}{Definition}[section]
 
\begin{definition} \label{def:userdata}
The user data include:
\begin{itemize}
\item Schedules;
\item Appointments;
\item User parameters.
\end{itemize}

\end{definition}
  
\begin{definition} \label{def:device}
A device is a PC, a Tablet or a Smartphone in which run the last version of his O.S.;
\end{definition}

\begin{definition} \label{def:schedule}
A Schedule is a set of time-ordered and not overlapping appointments where their starting times are fixed and they're linked to each other by a path covered with a specific transportation mean;
\end{definition}

\begin{definition} \label{def:bestTravelMean}
The bestTravelMean is the travel mean that satisfies the constraints on the appointment and on the schedule and that optimize the chosen criteria.
\end{definition}


\subsection{Acronyms}
%Acronyms used in the text:
\begin{itemize}
\item API: Application Programming Interface;
\item CRUD: Create/Read/Update/Delete;
\end{itemize}

%\reqref{req:R1}
%\goalref{goal:G1}

\section{Revision history}

\section{Reference Documents}
RASD document

\section{Document structure}
\begin{itemize}
\item Introduction: this section introduces the design document. It contains a justification of his utility
\item Architecture Design: this section is divided into different parts:
\begin{enumerate}
\item Overview: this section explains the division in tiers of our application and lists further details that were not specified in the RASD.
\item Components View: this sections gives a global view of the components of the application, other than a brief description, and describes how they communicate
\item Deployment view: this section shows the components that must be deployed to have the application running correctly
\item Runtime view: sequence diagrams are represented in this section to show how the components interact and collaborate during the fullfilling of a specific task. Each task is associated to each goal specified on the RASD document
\item Component interfaces: the interfaces between the components are presented in this section
\item Selected architectural styles and patterns : this section explains the architectural and design choices taken during the creation of the application
\end{enumerate}
\item Algorithms Design: this section describes the most important algorithmic parts. Pseudo code is used in order to hide unnecessary implementation details and abstract from a specific programming language
\item Requirements Traceability: this section aims to explain how the decisions taken in the RASD are linked to design elements
\end{itemize}



 

\chapter{Overall Description}

\section{Product Perspective}

\subsection{User Model}
A user is represented within the application by a set of parameters:
Some important informations about the user are held by the following ones:
\begin{itemize}
\item \textit{travelPass}: indicates if the user has a pass for public transportation;
\item \textit{hasBike};
\item \textit{hasCar}.
\end{itemize}

\subsection{Appointment Model}
An appointment is represented within the application by a set of parameters:
\begin{itemize}
\item \textit{duration}: the time extension of the appointment
\item \textit{startingTime} or \textit{timeInterval}: the first should be given if the starting hour is well-known (deterministic), otherwise a time interval in which the appointment will be held it's provided. 
%In particular: \begin{equation}timeInterval.start + duration \leq timeInterval.end \label{eqn:constrOnTimeInterval} \end{equation}
\item \textit{location}: identifies the coordinates of the place where the appointment will be held;
\item \textit{recurrent}: specifies if the appointment will be repeated over a fixed period of time;
\item \textit{peopleVariation}: represents a variation occourring when the user picks up or leaves off someone.
\end{itemize}

The life cycle of an appointment can be represented by the following statechart:
\begin{figure}[H]
\begin{center}
\includegraphics[width=\textwidth, keepaspectratio]{"images/stateChartAppointment".png}
\caption{Appointment statechart}
\end{center}
\end{figure}

A newly-created appointment is \textbf{Unscheduled}. It could remain \textbf{Unscheduled} either when edited or there isn't a possible arrangement when a schedule is performed. Otherwise it becomes \textbf{Scheduled} if there's a feasible way to arrange it. When a scheduled appointment is edited all the appointments in that schedule return to be \textbf{Unscheduled}, because it's possible that they bring to a different schedule. When a scheduled appointment is fullfilled it becomes \textbf{Archived} and stored in the schedule history. If this last one is a recurrent appointment it must be reinsert in the list of unscheduled appointment so it will become \textbf{Unscheduled} again. The user can cancel an appointment in every moment. 

\subsection{Schedule Model} \label{subsect:schmodel}
A schedule is a set of Appointments in a given day, ordered by the scheduler following the criteria described below.
A schedule is characterized by the following variables:
\begin{itemize}
\item \textit{date};
\item \textit{startingPosition}: is the starting location of the user (e.g. user's home);
\item \textit{startingNumberOfPeople}: the number of people that must reach the first appointment.
\item \textit{wakeUpTime}: it is the starting time from which the schedule should start arranging appointments.
\end{itemize}

\subsubsection{The optimization criteria} \label{subsubsect:optcriteria} 
The criteria that can be chosen for a schedule by the user  for the optimization are the following:
\begin{itemize}
\item \textit{Minimize carbon footprint}: the scheduler will try to minimize the amount of kilometers travelled in polluting means;
\item \textit{Minimize money spent}: the scheduler will try to avoid expensive means and to exploit the public ones (especially if the user has a pass) or going by bike or on foot;
\item \textit{Minimize travelling time}: the scheduler will compute the quickest possible path reaching all the appointments locations.
\end{itemize}

\subsection{Constraints}
Constraints are impositions on some parameters managed by the system during the process of scheduling the appointments. We can distinguish between constraints on schedule and contraints on the single appointment. These can be selected by the user when he inserts an appointment or when he requests a schedule, otherwise the constraints are initialized to default values. 

\subsubsection{Constraints on schedule} 
\begin{itemize}
\item \textit{Maximum travelling distance with a specific travel mean}: the user can set a maximum amount of km to travel with a travel mean;
\item\textit{ Travel means time slots}: user can specify a time interval in which a travel mean can be used;
\item \textit{User can deactivate a particular travel means};
\item \textit{User can select which travel means uses under certain weather condition}.
\end{itemize}

\subsubsection{Constraints on appointment}
\begin{itemize}
\item \textit{User can deactivate a particular travel means}.
\end{itemize}

\subsection{Class Diagram}
\begin{figure}[H]
\begin{center}
\includegraphics[width=\textwidth, keepaspectratio]{"images/classDiagram".png}
\caption{System Class Diagram}
\end{center}
\end{figure}

\section{Product Functions}
The following requirements are derived in order to fullfill the specified goals.
\\

Requirements for \goalref{goal:G1}:
\begin{enumerate}
\renewcommand\labelenumi{\textbf{R\theenumi}}
\item The system SPW to create a new appointment allowing user to specify all its parameters; \label{req:R1}
\item The system SPW to add constraints to a previously created appointment specifying which travel means have to be avoided during the travel to the meeting location; \label{req:R2}
\\

Requirements for \goalref{goal:G2}:

\item Allow the user to set the parameters of the schedule (\ref{subsect:schmodel}), or to accept the default values;
 
\item Allow the user to select the optimization criteria (\ref{subsubsect:optcriteria}) for the schedule;

\item The system S.B.A. to gather information from external APIs about: 
\begin{itemize}
\item travel options with related travel option data;
\item weather forecast;
\item strike days;
\end{itemize}

\item The system S.B.A. to select the best travel option according to the optimization criteria taking into account:
user constraint
travel option data 
weather forecast
information about strike day


%\item The system S.B.A. to compare different travel options with related travel option data according to the optimization criteria;

\item The system S.B.A. to avoid a particular travel mean under a certain weather condition;

\end{enumerate}

\subsection{ Assumptions, dependencies and constraints}

\begin{itemize}
\item The system should be able to retrieve information about public travel means. In particular:
\begin{itemize}
\item information about delays;
\item information on possible strike days.
\end{itemize}

\item The system should be able to retrieve information about shared means. In particular:
\begin{itemize}
\item position of the available ones;  
\item prices per time unit;
\end{itemize}

\item It's possible to retrieve all travel options and travel option data from external APIs;
\end{itemize}

%list of requirements given the goals



\chapter{Specific requirements}

\section{External Interface Requirements}
The application shows its best potential when running in a mobile device, for instance a smartphone or a tablet. This permits to extend the features and the automatic tasks of the application, thanks to the built-in device functionalities. However, a computer client version of the application can be installed, too.

\subsection{User interfaces}
The user can interact with the application through several graphical interfaces:
\begin{enumerate}
\item \textbf{Registration/login interface}: allows the user to insert credentials in order to registering or logging into the system;
\begin{figure}[H]
\begin{center}
\includegraphics[width=250pt, keepaspectratio]{"images/interfaces/login".jpg}
\caption{Registration/login interface}
\end{center}
\end{figure}

\item \textbf{User account interface}: user can specify his characteristics 
\begin{figure}[H]
\begin{center}
\includegraphics[width=250pt, keepaspectratio]{"images/interfaces/userprofile".jpg}
\caption{User account interface}
\end{center}
\end{figure}

\item \textbf{Home interface}: shows currently running schedule and displays some navigation links to other interfaces;
\begin{figure}[H]
\begin{center}
\includegraphics[width=250pt, keepaspectratio]{"images/interfaces/"home".jpg}
\caption{Home interface}
\end{center}
\end{figure}

\item \textbf{Appointment CRUD interface}: allows creating, showing and editing appointment parameters and related constraints;
\begin{figure}[H]
\begin{center}
\includegraphics[width=250pt, keepaspectratio]{"images/interfaces/appointment".jpg}
\caption{Appointment CRUD interface}
\end{center}
\end{figure}

\item \textbf{Appointments list interface}: provides a list of all inserted appointments, with the possibility to filter between non-scheduled/scheduled ones (includes the possibility to delete an item of the list);
\begin{figure}[H]
\begin{center}
\includegraphics[width=250pt, keepaspectratio]{"images/interfaces/"appointments".jpg}
\caption{Appointments list interface}
\end{center}
\end{figure}

\item \textbf{Schedules list interface}: display a list of the created schedules
\begin{figure}[H]
\begin{center}
\includegraphics[width=250pt, keepaspectratio]{"images/interfaces/"schedules".jpg}
\caption{Schedules list interface}
\end{center}
\end{figure}

\item \textbf{Schedule interface}: user can set parameters, contraints, optimization criteria and request a schedule creation for a given date;
\begin{figure}[H]
\begin{center}
\includegraphics[width=250pt, keepaspectratio]{"images/interfaces/"schedule".jpg}
\caption{Schedule interface}
\end{center}
\end{figure}

\item \textbf{Schedules result interface}: shows the computation of the requested schedules for a given date and asks the user to select one, then waits for confirmation for that;
\begin{figure}[H]
\begin{center}
\includegraphics[width=250pt, keepaspectratio]{"images/interfaces/"scheduleresults".jpg}
\caption{Schedules result interface}
\end{center}
\end{figure}

\item \textbf{Schedule progress interface}: permits to keep track of the completeness percentage, indicating the directions to be followed by the user in a map, in order to arrive to the next appointments;
\begin{figure}[H]
\begin{center}
\includegraphics[width=250pt, keepaspectratio]{"images/interfaces/schedule progress".jpg}
\caption{Schedule progress interface}
\end{center}
\end{figure}

\item \textbf{Tickets/rides reservation interface}: allows user to buy tickets for public travel means and/or reserve a ride for the shared travel means;
%\item \textbf{Appointments history interface}: shows a list of archived appointments;
\end{enumerate}

\subsubsection{Notes on User Interfaces}
The interfaces displayed before have been created through the IDE Xcode. In this way, if the application will be implemented, the GUI will look like exactly as shown above.

\subsection{Hardware interfaces}
Hardware interfaces are physical linking across which two or more separate components of a system exchange information. A hardware interface is described by the mechanical and electrical signals at the interface and the protocol for sequencing them. There are no interesting hardware interfaces in our scope.
%Our system relies on the following hardware interfaces:
%\begin{itemize}
%\item Mobile device: 
%\item Server: the subsystem is based on the client-server paradigm. 
%\item Travel means: gps on taxi and shared mean
%\end{itemize}

\subsection{Software interfaces}
Software interfaces are logical linking across which two or more separate applications running on a system exchange information. The most relevant software interface in our system is API. APIs are sets of subroutine definitions, protocols and clearly defined methods of communication, allowing data exchanging and service requests. There are several kinds of these:
\begin{itemize}
\item \textbf{Operating System APIs}: specify interface between applications and OS, permitting to access low level routines calls (for instance, to communicate with memory or with an internal device)
\item \textbf{Remote APIs}: DBMS expose a set of standards that the API user can adopt in order to manage the database data. SQL is the standard language for storing, manipulating and retrieving data in this context;
\item \textbf{Web API}: information can be exchanged through the internet by encapsulating it in HTTP request/response. Weather forecast, travel services, mapping systems offer this typology of API.
\end{itemize}

\subsection{Communications interfaces}
Communication interfaces allows two different architectures of the system to exchange information through communication channel. These non-homogeneous components of the system can communicate thanks to the following software interfaces and protocols:
\begin{itemize}
\item \textbf{Cellular connectivity}: mobile devices can connect to the internet thanks to LTE standard;
\item \textbf{GPS}: cellular can retrieve his coordinates position through NMEA protocol;
\item \textbf{QRCode}: associates a matrix of bits to an URL. QRCodes are present in most of the shared means, semplifying the booking of that.
%identify the nearby transportation
\end{itemize}


\section{Functional requirements}
\subsection{Scenarios}

Here are some scenarios that describe the usage of the system.

\subsubsection{Scenario 1} \label{scenario:1}
Luana has some problems in scheduling his daily appointments in fact she is always late. One of her friend told her that has been released a new application Travlendar+ that can be useful for scheduling appointments. Luana decides to download it and then she registers herself to the system by submitting his e-mail and a password. After the e-mail confirmation she inserts in the system her user parameters.


\subsubsection{Scenario 2} \label{scenario:2}
Giovanni will start the fourth year of his Master's degree. Surfing the internet, he finds out that his lesson schedule for the first semester has been published. Giovanni decides to fill in the application with his new appointments related to lessons attendance. In fact he knows where to go, at which time and day and for which amount of time. Since he knows that these events will going to happen for 3 months, he sets them as recurrent.

\subsubsection{Scenario 3} \label{scenario:3}
Edoardo wants to start training but he doesn't know what are the best hours in which he can run in accord to his appointments, he know only that he can run between 5 and 7 pm, for 45 minutes. he can insert this last appointment in the application whitout specify the exactly starting hour and the system will schedule it at best.

\subsubsection{Scenario 4} \label{scenario:4}
Federico has scheduled his appointments but at lunch time his son called him because he needed a ride for go back to home. Federico decided to help his son and so he brought him home. now the current running schedule is not more valid so he request to the system a reschedule of his appointment according to his position.  


\subsubsection{Scenario 5} \label{scenario:5}
John is ready to start is daily tasks in fact he has already scheduled his appointments through Travlenddar+ application. 15 minutes before the schedule starting time he recieves a notification by the application saying that there is a possible rearregment of his appointments by using shared travel means that can improve the optimization criteria selected by him. so he decides to accepts this new schedule.

\subsection{Use cases}

\subsubsection{User registration} \label{usecase:User registration}

\begin{longtable}{|p{14cm}|} \hline

\textbf{Name:} User registration \\ \hline

\textbf{Actors:} External User, external e-mail service \\ \hline

\textbf{Goals:} \goalref{goal:G1} \\ \hline

\textbf{Input Condition:} \\ \hline

\textbf{Event Flow:}

\begin{enumerate}
\item The user wants to register to the system, so he runs the application;
\item The system display the login/registration page;
\item The user fills the form (the form is present on the first page that the application display after the startup);
\item The user submit the filled form to the system;
\item The system send a confirmation e-mail to the user;
\item The system display a message in which the user is informed that he will recieve a confirmation e-mail;
\item The user confirm the registration by clicking a link in the received e-mail;
\item A confirmation message is sent to the application;
\item The user is redirected in his profile page inside the application;
\item The user specifies his parameters.
\end{enumerate}	\\ \hline

\textbf{Output Condition:} The registration is confirmed to the system; \\ \hline

\textbf{Exceptions:}

\begin{enumerate}
\item The e-mail given by the user is fake 
\item The user makes a typo during the insertion of his e-mail
\end{enumerate} \\ \hline

\textbf{Mapping on Requirements:}

\begin{itemize}
\item Events 1 through 3 are granted by \reqref{req:R1};
\item Evens 4 through 9 are granted by \reqref{req:R2};
\item Event 10 is granted by \reqref{req:R13}. 
\end{itemize}  \\ \hline

\end{longtable}

\begin{figure}[H]
\begin{center}
\includegraphics[width=400pt, keepaspectratio]{"images/RegistrationSequenceDiagram".png}
\caption{Registration sequence diagram}
\label{img:seqDiagrAppEditing00}
\end{center}
\end{figure}

\subsubsection{User log-in}

\begin{tabular}{|p{14cm}|} \hline

\textbf{Name:} User log-in \\ \hline

\textbf{Actors:} Registered User \\ \hline

\textbf{Goals}: \goalref{goal:G2}\\ \hline

\textbf{Input Condition:} The user is registered to the system \\ \hline

\textbf{Event Flow:} 

\begin{enumerate}
\item The user needs to log-in the application, so he runs it;
\item The system provides to the user a form to fill;
\item The user fills up the form with the his e-mail and his password; (as said in \ref{subsect:usermodel})
\item The user submits the form to the system;
\item The system checks the user identity and provides to the user the main application page .
\end{enumerate} \\ \hline

\textbf{Output Condition:} The user is logged-in to the system. \\ \hline

\textbf{Exceptions:} The user submits the form after having filled it with a wrong email or password. \\ \hline

\textbf{Mapping on requirements:}

\begin{itemize}
\item Events from 3 through 5 granted by \reqref{req:R5};
\item Event 6 grandet by \reqref{req:R4};
\end{itemize} \\ \hline

\end{tabular}

\includegraphics[width=400pt, keepaspectratio]{"images/LogInSequenceDiagram".png}
%\caption{Log-in sequence diagram}

\subsubsection{Recover credentials}\label{usecase:recovercredentials}
\begin{longtable}{|p{14cm}|} \hline
\textbf{Name:} Schedule appointments \\ \hline
\textbf{Actors:} Registered User, External Email Service \\ \hline
\textbf{Goals:} \goalref{goal:G3} \\ \hline
\textbf{Input Condition:} The user is registered to the system \\ \hline

\textbf{Event Flow:}
\begin{enumerate}
\item The user wants to recover his passoword;
\item The user requests to recover his passoword;
\item The system provides to the user the schedule form with his e-mail;
\item The user fills up the field of the form;
\item The user submits the form to the system;
\item The system sends the e-mail to the user with his password.
\end{enumerate}
\\ \hline

\textbf{Output Condition:} The user has recovered his password; \\ \hline

\textbf{Exceptions:}
\begin{enumerate}
\item The form it's left blank;
\item An invalid address is given;
\end{enumerate} \\ \hline

\textbf{Mapping on Requirements:}
\begin{itemize}
\item Actions 2 through 5 granted by requirement \reqref{req:R7}
\item Action 6 granted by requirement \reqref{req:R6}
\end{itemize}  \\ \hline

\end{longtable}

\begin{figure}[H]
\begin{center}
\includegraphics[width=450pt, keepaspectratio]{"images/PasswordRecoverySequenceDiagram".png}
\caption{Password recovery sequence diagram}
\label{img:seqDiagrPasswordRecovery}
\end{center}
\end{figure}

\subsubsection{Appointment creation} \label{usecase:appcreation}

\begin{longtable}{|p{14cm}|} \hline
\textbf{Name:} Appointment creation \\ \hline
\textbf{Actors:} Logged User \\ \hline
\textbf{Goals:} \goalref{goal:G4} \\ \hline
\textbf{Input Condition:} 
\begin{itemize}
\item The user is registered to the system 
\item The user is logged into the systems 
\end{itemize}
\\ \hline
\textbf{Event Flow:}
\begin{enumerate}
\item The user wants to add a new appointment to his schedule;
\item The user requests the appointments page
\item The system provides the appointments page
\item The user requests the creation of a new appointment to the application;
\item The system provides to the user a form to fill;
\item The user fills up the form with the parameters (specified in \ref{subsect:appointmentmodel}) and constraints (specified in \ref{subsubsect:constronappoint} about the new appointment;
\item The user submit the form to the system;
\item The system allocates the new appointment as Unscheduled (referring to statechart in figure; \label{fig:stchartApp})
\item The system sends a confirmation to the user.
\end{enumerate}	\\ \hline

\textbf{Output Condition:} The user has created a new appointment; \\ \hline

\textbf{Exceptions:}
\begin{enumerate}
\item Some fields of the form referring to parameters are left blank.
\end{enumerate} \\ \hline
\textbf{Mapping on Requirements:}
\begin{itemize}
\item Events 4 through 7 are granted by \reqref{req:R8}
\item Event 8 is granted by \reqref{req:R9}
\end{itemize}  \\ \hline


\end{longtable}

\begin{figure}[H]
\begin{center}
\includegraphics[width=400pt, keepaspectratio]{"images/sequenceDiagramAppointmentCreation".png}
\caption{Appointment creation sequence diagram}
\label{img:seqDiagrAppCreation}
\end{center}
\end{figure}

\subsubsection{Appointment editing}\label{usecase:appediting}
\begin{longtable}{|p{14cm}|} \hline
\textbf{Name:} Appointment editing \\ \hline
\textbf{Actors:} Logged User \\ \hline
\textbf{Goals:} (\goalref{goal:G5})\\ \hline
\textbf{Input Condition: The user is logged-in to the system} \\ \hline
\textbf{Event Flow:}
\begin{enumerate}
\item The user wants to modify an appointment of his schedule;
\item The user selects the appointment to modify;
\item The system provides to the user the appointment form with all the parameters and constraints;
that were specified yet by the user; 
\item The user edits the fields of the form;
\item The user submit the form to the system;
\item The system set the appointment as Unscheduled with the new parameters (referring to statechart in figure \ref{fig:stchartApp});
\item The system sends a confirmation to the user.
\end{enumerate}	\\ \hline
\textbf{Output Condition:}  The user has modified an appointment; \\ \hline
\textbf{Exceptions:}
\begin{enumerate}
\item Some fields of the form referring to parameters are left blank.
\end{enumerate} \\ \hline
\textbf{Mapping on Requirements:}
\begin{itemize}
\item Events 3 through 5 are granted by \reqref{req:R10}
\item Event 6 is granted by \reqref{req:R11}
\end{itemize}  \\ \hline

\end{longtable}

\begin{figure}[H]
\begin{center}
\includegraphics[width=400pt, keepaspectratio]{"images/seqDiagramAppointmentEditing".png}
\caption{Appointment editing sequence diagram}
\label{img:seqDiagrAppEditing00}
\end{center}
\end{figure}

\subsubsection{Schedule appointments}\label{usecase:scheduleappointments}
\begin{longtable}{|p{14cm}|} \hline
\textbf{Name:} Schedule appointments \\ \hline
\textbf{Actors:} Logged User, External API \\ \hline
\textbf{Goals:} \goalref{goal:G6} \\ \hline
\textbf{Input Condition:} The user is logged-in to the system \\ \hline

\textbf{Event Flow:}
\begin{enumerate}
\item The user wants to schedule his appointments;
\item The user requests the creation of a new schedule for the current day;
\item The system provides to the user the schedule form with all the parameters, optimization criteria and the constraints;
\item The user fills up the fields of the form;
\item The user submits the form to the system;
\item The system retrieves information from external APIs about; travel options and related travel option data, weather forecast and strike days;
\item The system retrieves about the Travel Option Data of the newly created Schedule
\item The system stores the Schedule, together with his travel option data and stores it as Saved (\ref{fig:stchartApp}) with the appointment selected by the user.
\end{enumerate}	\\ \hline

\textbf{Output Condition:} The user has created a valid schedule of his appointments; \\ \hline

\textbf{Exceptions:}
\begin{enumerate}
\item Some fields of the form referring to schedule variables and optimization criteria are left blank. The parameters of schedule constraints could also be left blank since they will assume default values;
\item It's not possible to list the appointments as a Valid Schedule, so the schedule is Discarded (\ref{fig:stchartApp})
\end{enumerate} \\ \hline
\textbf{Mapping on Requirements:}
\begin{itemize}
\item Events 3 through 5 are granted by \reqref{req:R12} through \reqref{req:R14};
\item Event 6 and 7 is granted by \reqref{req:R15};
\item Event 8 is granted by \reqref{req:R16} and \reqref{req:R17}.
\end{itemize}  \\ \hline

\end{longtable}

\begin{figure}[H]
\begin{center}
\includegraphics[width=450pt, keepaspectratio]{"images/SequenceDiagramSchedule".png}
\caption{Schedule appointents sequence diagram}

\label{img:seqDiagrAppEditing00}
\end{center}
\end{figure}

\subsubsection{Schedule selection}
\begin{longtable}{|p{14cm}|} \hline
\textbf{Name:} Multiple Schedules creation \\ \hline
\textbf{Actors:} Logged User \\ \hline
\textbf{Goals:} \goalref{goal:G7} \\ \hline
\textbf{Input Condition:} 
\begin{itemize}
\item The user is registered to the system; 
\item The user is logged in to the systems. 
\end{itemize}
\\ \hline
\textbf{Event Flow:}
\begin{enumerate}
\item The user wants to compare multiple schedules;
\item The user requests the schedules page;
\item The system provides the schedules page;
\item The user selects a schedule to be run;
\item The system display the mainpage with the schedule results (\ref{def:schedulingResult}).
\end{enumerate}	\\ \hline
\textbf{Output Condition:} The user selects a schedule to be run \\ \hline
\textbf{Exceptions:}
\\ \hline
\textbf{Mapping on Requirements:}
\begin{itemize}
\item Events are granted by the requirment \reqref{req:R18}
\end{itemize}  \\ \hline

\end{longtable}
\label{usecase:ScheduleSelection}

\begin{figure}[H]
\begin{center}
\includegraphics[width=400pt, keepaspectratio]{"images/ScheduleSelectionSequenceDiagram".png}
\caption{Schedule selection sequence diagram}
\label{img:ScheduleSelection}
\end{center}
\end{figure}


\subsubsection{Booking phase} \label{usecase:Booking Phase}

\begin{longtable}{|p{14cm}|} \hline

\textbf{Name:}Booking phase \\ \hline

\textbf{Actors:} Logged User, External APIs \\ \hline

\textbf{Goals:} \goalref{goal:G8} \\ \hline

\textbf{Input Condition:}

\begin{itemize}
\item The user must be logged in to the system;
\item The user must have selected a schedule to be run;
\item The user must have linked to the system his external accounts;
\item The user would like to buy the tickets for the travel means involved in the running schedule.
\end{itemize}
 \\ \hline

\textbf{Event Flow:}

\begin{enumerate}
\item The System, after the user have selected a schedule, asks to the user if he want to buy the ticket for the running schedule;
\item The User confirm to the system his intention;
\item The System perform a call to the travel means APIs for buying the ticket;
\item The APIs send back a confirmation message of the purchase;
\item The system send a confirmation message to the user.
\end{enumerate}	\\ \hline

\textbf{Output Condition:} The User recieve the confirmation message; \\ \hline

\textbf{Exceptions:}

\begin{enumerate}
\item The user doesn't have enough money in his card to complete the transaction;
\item there aren't free sits in one of the selected travel means; 
\end{enumerate} \\ \hline

\textbf{Mapping on Requirements:}

%\begin{itemize}
Events 3 through 5 granted by \reqref{req:R19}; \\ \hline
%\end{itemize}  \\ \hline

\end{longtable}

\begin{figure}[H]
\begin{center}
\includegraphics[width=400pt, keepaspectratio]{"images/BookingPhaseSequenceDiagram".png}
\caption{Booking phase sequence diagram}
\label{img:seqDiagrAppEditing00}
\end{center}
\end{figure}


\subsubsection{Dynamic directions} \label{usecase:Dynamic Directions}

\begin{longtable}{|p{14cm}|} \hline

\textbf{Name:}Dynamic Directions\\ \hline

\textbf{Actors:} Logged User, External APIs, GPS \\ \hline

\textbf{Goals:} \goalref{goal:G9} \\ \hline

\textbf{Input Condition:}

\begin{itemize}
\item The user must be logged in to the system;
\item The user must have a running schedule;
\end{itemize}
 \\ \hline

\textbf{Event Flow:}

\begin{enumerate}
\item The user requests the Directions for the travel to the system;
\item The system retrieves the user position from his GPS;
\item The system retrives from external APIs the directions to give to the user based on his position;
\item The system display to the user the updated map and the directions that him must follow in order to arrive to the next appointment
\item User doesn't need more directions so he closes the dynamic map.
\end{enumerate}	\\ \hline

\textbf{Output Condition:} The User is satisfied with the information gathered until this moment so he decides to close the dynamic map; \\ \hline

%\textbf{Exceptions:}
%
%\begin{enumerate}
%\item 
%\end{enumerate} \\ \hline

\textbf{Mapping on Requirements:}

\begin{itemize}
	\item Events 2 granted by \reqref{req:R20}
	\item Event 3 granted by \reqref{req:R21}
	\item Event 4 granted by \reqref{req:R22}
\end{itemize}  \\ \hline

\end{longtable}

\begin{figure}[H]
\begin{center}
\includegraphics[width=400pt, keepaspectratio]{"images/DynamicDirectionsSequenceDiagram".png}
\caption{Dynamic directions sequence diagram}
\end{center}
\end{figure}

In this sequence diagram the loop of the actions described in the rectangle continues until the dynamicMap is closed (closeDynamicMap).


\subsubsection{Notify Shared Means}\label{usecase:Notify Shared Means}
\begin{longtable}{|p{14cm}|} \hline
\textbf{Name:} Notify Shared Means \\ \hline
\textbf{Actors:} Registered User, External Api \\ \hline
\textbf{Goals:} \goalref{goal:G10} \\ \hline

\textbf{Event Flow:}
\begin{enumerate}
\item The system requests information to an external API about Shared Travel Means;
\item The external API service respond to the system with the information requested;
\item The system with the gathered information computes if there is a better path for the user according to the chosen constraints and optimization criteria;
\item if the path is found by the system is sent a notification to the user;
\end{enumerate}
\\ \hline

\textbf{Output Condition:} a better path is found; \\ \hline

\textbf{Exceptions:}
%\begin{enumerate}
%\end{enumerate} \\ \hline

\textbf{Mapping on Requirements:}
\begin{itemize}
\item Actions 1 and 2 are granted by \reqref{req:R24}
\item Actions 4 is granted by \reqref{req:R23}
\end{itemize}  \\ \hline

\end{longtable}

\begin{figure}[H]
\begin{center}
\includegraphics[width=400pt, keepaspectratio]{"images/NotifySharedMeansSequenceDiagram".png}
\caption{Notify Shared Means sequence diagram}
\end{center}
\end{figure}

\subsection{Use Case Diagram}

\begin{figure}[H]
\begin{center}
\includegraphics[width=450pt, keepaspectratio]{"images/useCaseDiagram".png}
\caption{Use Case Diagram}
\label{img:seqDiagrPasswordRecovery}
\end{center}
\end{figure}

\subsection{Notes on diagrams}
In these diagrams we assume that the user has ran the application. In particular, if the actor is a \textit{Registered User} or an \textit{Enternal User} then we assume that the user is facing the login page. On the other hand, if the actor is a \textit{Logged User} then is assumed that it is on the home page.

\section{Performance Requirements} \label{\sec:performanceRequirements}
The user must be notified in real time when a shared travel mean can be booked in order to provide a better mobility option,
since these kind of transportation can remain available for a limited amount of time.

Other performance requirements can't be easily expressed because they depends heavily on external services and on the device in which the application is run. For example the time needed to create a schedule is influenced by the promptness of the APIs.

Anyway an upper bound of 5 seconds for the creation of a schedule is given. Moreover, the position of the user during the progress of a schedule must be track with a maximum delay of 100ms.

\section{Design Constraints}

\subsection{Standard compliance}
Our system conforms to OAuth2 \footnote{Open Authentication, an industry-standard protocol for authorization. It focuses on client developer simplicity while providing specific authorization flows for web applications, desktop applications, mobile phones, and living room devices.} to handle the registration and login process. 
Moreover, HTTPS \footnote{HTTPS is a communications protocol for secure communication over a computer network. The main motivation for HTTPS is authentication of the visited website and protection of the privacy and integrity of the exchanged data.} protocol is used to guarantee secure treatment of user's sensitive data.

\subsection{Hardware limitations}
The bottleneck on the performance of the system is represented by the network infrastructure capability. In particular the most affected activity is the schedule computation.
\subsubsection{Analysis}
We can assume that the upload and download speeds of the APIs server is respectively 50Mbit/s and 150Mbit/s per client and that a request and a response weights are 20KB and 180KB.
We can consider 2 cases:
\begin{itemize}
\item 10Mbit/s and 5Mbit/s download and upload speeds respectively of the client;
\item 100Mbit/s and 50Mbit/s download and upload speeds respectively of the client.
\end{itemize}

\begin{table}[htbp]
  \centering
    \begin{tabular}{|r|r|r|r|r|}
    \hline
    \multicolumn{1}{|l|}{\textbf{No. Appointments}} & \multicolumn{1}{l|}{\textbf{No. APIs calls}} & \multicolumn{1}{l|}{\textbf{Request [MB]}} & \multicolumn{1}{l|}{\textbf{Response [MB]}} & \multicolumn{1}{l|}{\textbf{Tot [sec]}} \\ \hline
    %\midrule
    5 (2)    & 4   & 0.625 & 5.625 & 0.687 \\ \hline
    %\midrule
    10 (4)   & 28 & 4.375 & 39.375 & 4.812 \\ \hline
    %\midrule
    15 (7)   & 5047 & 788.6 & 7097.3 & 867.453 \\ \hline
    %\bottomrule
    \end{tabular}%
    \caption{Total times in the case of 10Mbit/s download speed and 5Mbit/s upload speed}
  \label{tab:addlabel}%
\end{table}%

\begin{table}[htbp]
  \centering
    \begin{tabular}{|r|r|r|r|r|}
    \hline
    \multicolumn{1}{|l|}{\textbf{No. Appointments}} & \multicolumn{1}{l|}{\textbf{No. APIs calls}} & \multicolumn{1}{l|}{\textbf{Request [MB]}} & \multicolumn{1}{l|}{\textbf{Response [MB]}} & \multicolumn{1}{l|}{\textbf{Tot [sec]}} \\ \hline
    %\midrule
    5 (2)    & 4   & 0.625 & 5.625 & 0.125 \\ \hline
    %\midrule
    10 (4)   & 28 & 4.375 & 39.375 & 0.875 \\ \hline
    %\midrule
    15 (7)   & 5047 & 788.6 & 7097.3 & 157.719 \\ \hline
    %\bottomrule
    \end{tabular}%
    \caption{Total times in the case of 100Mbit/s download speed and 50Mbit/s upload speed.}
  \label{tab:addlabel}%
\end{table}%

In the first column the numbers between brackets rempresent the number $n$ of appointments with variable starting time, so that can be arranged differently relative to each other, changing their order in the schedule. Then the number of calls to External APIs that should be done is calculated, in case of a brute-force approach in the scheduling algorithm. Therefore the number of calls is proportional to $n!$.
Finally the total amount of time is calculated considering the previous assumptions.

We can realize that the number of calls to external APIs should be minimized in order to fullfill the requirement on performance expressed in \ref{\sec:performanceRequirements}.

\section{Software System Attributes}

\subsection{Reliability}
The system should guarantee that from the data retrieved is always constructed the most convenient valid schedules, according to user preferences and constraints, if it exists.

\subsection{Availability}
The system should be accessible 24 hours per day and should be available 99,9\% of the time (up to 8,76 hours per year of downtime). Anyway the availability of the features involving the use of external services can't be directly controlled. In particular the availability of the feature $j$ is given by:

\begin{equation}
A_j = a_0 \prod_{i=1}^n a_i 
\end{equation} 

where each $a_i$ represent the availability of the external service $i$ used and $a_0$ is the availability of the application. For instance, in the case of a schedule creation:

\begin{figure}[H]
\begin{center}
\includegraphics[width=\textwidth, keepaspectratio]{"images/availability diagram".jpg}
\caption{A failure in one of the chain of request to the APIs causes the entire process to break down}
\end{center}
\end{figure}

\subsection{Security}
The identity of the user must be verified through a login phase.
User's characteristics must be protected during transmission from client to server throughout the registration. User credentials are cryptographied and then saved.
\subsection{Maintainability}
The system should be open to modifications.
In particular the application should be able to consider new travel means, new scheduling optimization criteria and new constraints.
Moreover, also the GUI should be easily editable, so that can adapt to new operating systems.

\subsection{Portability}
The system should be adaptable to run in all the devices (\textbf{link alla definizione}) considered.

\chapter{Alloy analysis}

\begin{lstlisting}
module Travlendar

open util/integer as Integer
open util/boolean as boolean
// impose ordering among relations
open util/ordering[Date] as DO 
open util/ordering[Time] as TO
open util/ordering[ScheduledAppointment] as SAO

-----------------------------------------------
sig Date{}
-----------------------------------------------
sig Time{}

-----------------------------------------------
sig TimeSlot{
	start: one Time,
	end: one Time
}

-----------------------------------------------
sig OptimizingCriteria{}{
	this in Schedule.optimizingCriteria
}

-----------------------------------------------
sig User{
	hasCar: one Bool,
	hasBike: one Bool,
	hasPass: one Bool
}{
	this in Appointment.user
}

-----------------------------------------------
sig Schedule{
	date: one Date,
	wakeUpTime: one Time,
	optimizingCriteria: one OptimizingCriteria,
	constraints : set ConstraintOnSchedule,
	initialNumberOfPeopleInvolved : one Int
}

// there aren't schedules that doesn't have at least one ScheduledAppointment
fact NoScheduleUnlinked{
	all s:Schedule | s in ScheduledAppointment.schedule  
}

// we can't have two different schedule for the same day with the same optimizing criteria
fact DifferentOptimizingCriteriaForScheduleOfTheSameDay{
	all s,s1 : Schedule| s != s1 and s.date = s1.date  => s.optimizingCriteria != s1.optimizingCriteria 
}

// no schedule of scheduled appointment belonging to different users
fact belongingCoeherence{
	no sa1,sa2 :ScheduledAppointment | sa1.schedule = sa2.schedule 
	and sa1.appointment.user != sa2.appointment.user
}

-----------------------------------------------
sig Appointment{
	user: one User,
	date: one Date,
	startingTime: lone Time,
	timeSlot: lone TimeSlot,
	constraints: set ConstraintOnAppointment,
	variationNumberInvolvedPeople: one Int
}
{	
	// startingTime and timeSlot are mutually exclusive
	startingTime = none => timeSlot != none
	timeSlot = none => startingTime != none
	startingTime in TO/nexts[timeSlot.start] and startingTime in TO/prevs[timeSlot.end]
}

// all appointments must be scheduled into some ScheduledAppointment
fact AppointmentAssociationCoherence {
	all a : Appointment | a in ScheduledAppointment.appointment
}

-----------------------------------------------
sig ScheduledAppointment{
	schedule: one Schedule,
	date: one Date,
	appointment: one Appointment,
	startingTravelTime: one Time,
	ETA: one Time,
	endingTime: one Time,
	numberOfInvolvedPeople: one Int,
	weather: one Weather
}
{	startingTravelTime in TO/prevs[endingTime]
	// ETA must be between starting and ending times
	ETA in TO/nexts[startingTravelTime] and ETA in TO/prevs[endingTime] 
	date = schedule.date																				
	date = appointment.date																			
	numberOfInvolvedPeople >= 0
	startingTravelTime in TO/nexts[schedule.wakeUpTime]
}

// appointments starting time ordering coherence
fact ScheduledAppointmentsOrderingConsistence {
	all a1,a2 : ScheduledAppointment | a1.schedule = a2.schedule and a1 in SAO/prevs[a2] 
	=> a1.startingTravelTime in TO/prevs[a2.startingTravelTime]
}

// if two scheduled appointments are relative to the same appointment then they must 
//belong to different schedules
fact AppointmentOnMultipleSchedules {
	all s1,s2 : ScheduledAppointment | s1 != s2 and s1.appointment = s2.appointment 
	=> s1.schedule != s2.schedule
}

// Starting time of a scheduled appointment must coincide with the starting time of the 
//non-scheduled appointment if timeSlot is not specified
fact StartingTimeCoherence{
	all s : ScheduledAppointment | s.appointment.timeSlot = none =>
	s.startingTravelTime = s.appointment.startingTime
}

// The number of involved people in each travel is coherent with the number of seats of each 
//travel mean used for it
fact numberOfPeopleInvolvedCoherentWithSeats{
	no p : Path | p.travelMean.seats < p.source.numberOfInvolvedPeople
}

// The selected travel mean for a path can provide enough seats for the people involved in the appointment
fact coherenceOnNumberOfInvolvedPeople{
	all sa:ScheduledAppointment, a:Appointment, s:Schedule | 
	sa in schedule.s and a in sa.appointment =>
	sa.numberOfInvolvedPeople = add[s.initialNumberOfPeopleInvolved,
		sum e : SAO/prevs[sa] | e.appointment.variationNumberInvolvedPeople]
}

----------------------------------------------------
sig Path{
	lenght: one Int,
	source: one ScheduledAppointment,
	dest: one ScheduledAppointment,
	travelMean: one TravelMean
}
{
	lenght >= 0
	source != dest
	source.schedule = dest.schedule
}

// No path linking to scheduled appointment belonging to different user
fact SchedulePathBelongingConsistency{ 
	all p:Path | p.source.appointment.user = p.dest.appointment.user
}

// No scheduled appointment unreachable
fact AllScheduledAppointmentCanBeReachedByPaths{
	all s : ScheduledAppointment | s in (Path.source + Path.dest)
}

// No path belonging to a user with a travel mean that he/she doesn't own
fact TravelMeanConsistency{
	no p:Path| (p.travelMean = Bike and p.source.appointment.user.hasBike in False) or 
	(p.travelMean = Car and p.source.appointment.user.hasCar in False) 
}

// No path linking two scheduled appointments that are not sequential
fact PathOrderConsistency{
	all p:Path | p.source = SAO/prev[p.dest] 
} 

-----------------------------------------------------
abstract sig Constraint{
	travelMean: one TravelMean,
	maxTravelDistance: one Int // with travel mean above
}
{maxTravelDistance >= 0} //if 0 implies that the travel mean is deactivated

sig ConstraintOnAppointment extends Constraint{
}
{maxTravelDistance = 0}

sig ConstraintOnSchedule extends Constraint{
	weather: set Weather, 
	timeSlot: lone TimeSlot,	 
	strikeDate: lone Bool
	// in which travel mean can't be used
}{
	weather != none => maxTravelDistance = 0
	weather != none => (timeSlot .start = timeSlot.end)
}

fact NoConstraintUnlinked{
	all c : Constraint | c in Appointment.constraints or c in Schedule.constraints
}

-----------------------------------------------
abstract sig Weather{} 

one sig  Sunny extends Weather{}
one sig Snowy extends Weather{}
one sig Rainy extends Weather{}
one sig Foggy extends Weather{}
one sig Cloudy extends Weather{}

fact NoWeatherUnlinked {
	all w : Weather | w in ConstraintOnSchedule.weather 
}

------------------------------------------------
abstract sig TravelMean{ 
	seats: one Int
}
//100 seats encode an unbounded number of seats

abstract sig PublicTravelMean extends TravelMean{}
abstract sig SharedTravelMean extends TravelMean{}
abstract sig PrivateTravelMean extends TravelMean{}

sig Train extends PublicTravelMean{}{
	seats=100
}
sig Bus extends PublicTravelMean{}{
	seats=100
}
sig Tram extends PublicTravelMean{}{
	seats=100
}
sig Taxi extends PublicTravelMean{}{
	seats=3
}
sig Underground extends PublicTravelMean{}{
	seats=100
}
sig CarSharing extends SharedTravelMean{}{
	seats=5
}
sig BikeSharing extends SharedTravelMean{}{
	seats=1
}
sig Bike extends PrivateTravelMean{}{
	seats=1
}
sig Car extends PrivateTravelMean{}{
	seats=5
}
sig Walking extends PrivateTravelMean{}{
	seats=100
}

-----------------------------------------------
// if a schedule appointment is on a schedule, it must be in the list of scheduled appointments of that schedule
ScheduledAppointmentAndScheduleBiunivocity : check{
	all a : ScheduledAppointment | all s : Schedule | a.schedule = s => a in schedule.s
} for 8
-----------------------------------------------------
// if a scheduled appointment is in a schedule, they must have the same date
EveryAppoinmentOfScheduleIsInItsDay : check {
	all s : Schedule, a : ScheduledAppointment | a.schedule = s => s.date = a.date
}
for 8
-----------------------------------------------------
ScheduleIsOwnedByOnlyOneUser : check {
	all s1,s2:ScheduledAppointment | s1.schedule = s2.schedule implies s1.appointment.user = s2.appointment.user
} for 8
-----------------------------------------------------
// all scheduled appointment date must be equal to its original appointment date
AppointmentDateIsEqualToItsScheduledAppt : check {
	all a : ScheduledAppointment | a.date = a.appointment.date
} for 8
-----------------------------------------------------
//source and dest of each path must refer to the same schedule
sourceDestOfPathCoherence : check {
	all p:Path | p.dest.schedule = p.source.schedule
} for 8
-----------------------------------------------------
// scheduled appointments are ordered according to their path precedence
PathOrderConsistency : check {
	all p:Path | p.dest = SAO/next[p.source]
} for 8
-----------------------------------------------------
/* if two appointments belong to different user then they should belong to different 
    schedule since a schedule belongs to only one user */ 
TwoAppointmentsWithSameUserNoInSameSchedule : check {
	all a,a1:Appointment | a.user != a1.user => appointment.a.schedule != appointment.a1.schedule
} for 8
-----------------------------------------------------

// return all the paths of the specified schedule
fun pathOfSchedule [s : Schedule] : set Path {
	(source+dest).(schedule.s)
}

// return the total distance travelled in the specified schedule with a specific travel mean
fun distanceTravelledWithMeanInSchedule [s : Schedule, t: TravelMean] : Int {
	sum e : (pathOfSchedule[s]) & (travelMean.t) | e.lenght
}

// checks if the specified path links two scheduled appointment in the given schedule
pred doesPathBelongToSchedule [s:Schedule, p:Path] {
	some sa:ScheduledAppointment | sa in (Path.source + Path.dest) and sa.schedule=s
}

// checks if a schedule satisfying the specified constraint exists
pred doesConstraintSatisfySchedule (c : ConstraintOnSchedule){
	some s : Schedule |  all p:Path |
	doesPathBelongToSchedule[s, p] and
	p.travelMean  != c.travelMean or (p.travelMean = c.travelMean 
					and (schedule.s).weather not in c.weather 
					and ((c.timeSlot.start = c.timeSlot.end) 
						or ((schedule.s).ETA in TO/prevs[c.timeSlot.start]
						or (schedule.s).startingTravelTime in TO/nexts[c.timeSlot.end] ))
					and (c.strikeDate in True implies p.travelMean not in PublicTravelMean)
					and (distanceTravelledWithMeanInSchedule[s,c.travelMean] < c.maxTravelDistance )
				)
}

// checks if a scheduled appointment satisfying the specified constraint exists
pred doesConstraintSatisfiesAppointment(c : ConstraintOnAppointment){
	some s: Schedule | all sa:ScheduledAppointment, a: Appointment, p: Path |
		sa in schedule.s and sa.appointment = a and p.dest=sa
		=> p.travelMean != c.travelMean
}

pred NoOverlappingScheduledAppointmentInSchedule (s : Schedule) {
	all s1,s2 : ScheduledAppointment |  s1.schedule = s and  s2.schedule = s and s1 != s2 
	=> s1.endingTime in TO/prevs[s2.startingTravelTime] - s2.startingTravelTime
	or s2.endingTime in TO/prevs[s1.startingTravelTime] - s1.startingTravelTime
}

// checks if a valid schedule exists (see def. on RASD document)
pred validSchedule {
	some s: Schedule | all cs:ConstraintOnSchedule, ca:ConstraintOnAppointment  | 
	cs in s.constraints and ca in (schedule.s).appointment.constraints
	and doesConstraintSatisfySchedule[cs] and doesConstraintSatisfiesAppointment[ca]
	and NoOverlappingScheduledAppointmentInSchedule[s]
}
-----------------------------------------------------

// new appointment insertion action
pred InsertNewAppointment[u,u1 : User, a1 : Appointment]{
	//precondition
	all a : Appointment | a.user = u and a1.user != a.user
	//postcondition
	user.u1 = user.u + a1
}

// checks if the insertion of a new appointment is correct
InsertNewAppointmentIsCorrect : check {
	all u,u1 : User, a,a1 : Appointment | a.user = u and a1.user != a.user and InsertNewAppointment[u,u1,a1] 
	=> a1 not in user.u and a1 in user.u1
} for 8
-----------------------------------------------------

// editing appointment action
pred ModifyAppointment[a,a1 : Appointment, u : User]{
	//precondition
	a in user.u and a1 not in user.u
	//postcondition
	a1 in user.u and a not in user.u
}

// checks if the modification of a new appointment is correct
ModifyAppointmentIsCorrect : check {
	all u : User, a,a1 : Appointment | a in user.u and a1 not in user.u and ModifyAppointment[a,a1,u]
	=> a1 in user.u and a not in user.u
} for 8


pred show(){}

run { show and validSchedule} for 7 but 4 Int

\end{lstlisting}


%%APPENDICI==============================================
\appendix
%%\input{appendiceA}

\end{document}